\section*{Chapitre \ref{ch:g9:kinetic-energy-safety} --- L'énergie cinétique et la sécurité routière}

\begin{solution}{exo:g9:kinetic-energy-safety:1}
Le \omterm{def:g9:kinetic-energy-safety:joule}{joule} (\unit{J})~: quelques \omterm{def:g9:kinetic-energy-safety:joule}{joules} montent ce livre sur son
étagère~; le cœur en dépense environ un par battement.
$E_k = \frac12 m v^2$ --- $m$ en \omterm{def:g3:measuring-mass:units}{kilogrammes}, $v$ en \omterm{def:g2:measuring-length:units}{mètres} par
seconde, $E_k$ en \omterm{def:g9:kinetic-energy-safety:joule}{joules}.
\end{solution}

\begin{solution}{exo:g9:kinetic-energy-safety:2}
Lièvre~: $0.5 \times 2 \times 100 = \qty{100}{J}$. Joueur~:
$0.5 \times 80 \times 25 = \qty{1000}{J}$.
\end{solution}

\begin{solution}{exo:g9:kinetic-energy-safety:3}
Une moitié de \omterm{def:g3:measuring-mass:mass}{masse} en plus~: facteur $1.5$. Une moitié de \omterm{def:g7:motion-average-speed:formula}{vitesse} en
plus~: facteur $1.5^2 = 2.25$. Le réglage de la \omterm{def:g7:motion-average-speed:formula}{vitesse} l'emporte ---
il l'emporte toujours, puisqu'il porte le carré.
\end{solution}

\begin{solution}{exo:g9:kinetic-energy-safety:4}
La \omterm{prop:g9:kinetic-energy-safety:stopping}{distance de réaction} --- le trajet d'une seconde, croissant
proportionnellement à $v$~; la \omterm{prop:g9:kinetic-energy-safety:stopping}{distance de freinage} --- la mise à la
retraite de $E_k$, croissant comme $v^2$.
\end{solution}

\begin{solution}{exo:g9:kinetic-energy-safety:5}
À \qty{50}{km/h}~: $13$ sur \qty{27}{m} --- environ la moitié. À
\qty{130}{km/h}~: $85$ sur \qty{121}{m} --- quelque soixante-dix pour
cent. La composante au carré dépasse la composante proportionnelle~:
à grande \omterm{def:g7:motion-average-speed:formula}{vitesse}, le freinage dévore le total.
\end{solution}

\begin{solution}{exo:g9:kinetic-energy-safety:6}
En \omterm{def:g5:heat-and-insulation:heat}{chaleur}, dans les disques de frein --- le rougeoiement d'ambre
après une descente de montagne, l'odeur de frein chaud. L'\omterm{def:g5:everyday-energy:energy}{énergie} se
transmet, elle ne s'efface jamais~: la proposition de la chaîne du
chapitre d'\omterm{def:g5:everyday-energy:energy}{énergie} de l'enfance, toujours aux commandes.
\end{solution}

\begin{solution}{exo:g9:kinetic-energy-safety:7}
$36 \div 3.6 = \qty{10}{m}$~; puis \qty{20}{m}~; puis \qty{30}{m} ---
en marche proportionnelle à la \omterm{def:g7:motion-average-speed:formula}{vitesse}, comme le doit une portion
d'une seconde.
\end{solution}

\begin{solution}{exo:g9:kinetic-energy-safety:8}
Ton $E_k$ est fixée par ta \omterm{def:g3:measuring-mass:mass}{masse} et par la \omterm{def:g7:motion-average-speed:formula}{vitesse} de la voiture ---
aucune sangle ne la change. La ceinture civilise la \emph{retraite}~:
elle arrête le corps sur une distance et une \omterm{def:g2:measuring-time:duration}{durée} plus longues,
doucement, au lieu de le laisser s'arrêter contre le tableau de bord,
d'un seul coup.
\end{solution}

\begin{solution}{exo:g9:kinetic-energy-safety:9}
Le freinage se met à l'échelle du carré de la \omterm{def:g7:motion-average-speed:formula}{vitesse}~: à
\qty{60}{km/h} (le double), $8 \times 4 = \qty{32}{m}$~; à
\qty{90}{km/h} (le triple), $8 \times 9 = \qty{72}{m}$.
\end{solution}

\begin{solution}{exo:g9:kinetic-energy-safety:10}
À \qty{30}{km/h}~: réaction $\approx \qty{8.3}{m}$, freinage
$13 \times (30/50)^2 \approx \qty{4.7}{m}$ --- arrêt en environ
\qty{13}{m}, contre \qty{27}{m} à cinquante~: moins de la moitié de
la route. \omterm{def:g5:everyday-energy:energy}{Énergies}~: $(30/50)^2 = 0.36$ --- à peine le tiers de
l'\omterm{def:g5:everyday-energy:energy}{énergie} de choc. Les deux convainquent~; la ligne prête pour
l'affiche est d'ordinaire celle de la route --- «~à 30, tu t'arrêtes
avant l'enfant~; à 50, tu l'atteins encore lancé~».
\end{solution}

\begin{solution}{exo:g9:kinetic-energy-safety:11}
Ne doubler que la part du freinage --- la pluie allonge la retraite,
non le cerveau~: à \qty{50}{km/h}, $14 + 26 = \qty{40}{m}$~; à
\qty{90}{km/h}, $25 + 80 = \qty{105}{m}$. La seconde de réaction est
la même sur route sèche et mouillée.
\end{solution}

\begin{solution}{exo:g9:kinetic-energy-safety:12}
Camion~: $v = 25$ \unit{m/s}~;
$E_k = 0.5 \times 20000 \times 625 = \qty{6.25e6}{J}$. Voiture qui
égale~: $0.5 \times 1000 \times v^2 = \num{6.25e6}$ donne
$v^2 = 12500$, $v \approx \qty{112}{m/s} \approx \qty{400}{km/h}$ ---
aucune voiture de route n'en approche. Un camion à \omterm{def:g7:motion-average-speed:formula}{vitesse}
d'autoroute porte des \omterm{def:g5:everyday-energy:energy}{énergies} de voiture de course avec des freins de
\omterm{def:g4:weight-and-mass:weight}{poids} lourd~: quand ils surchauffent et lâchent dans une longue
descente, seul un lit de gravier peut mettre des mégajoules à la
retraite sans dégâts. Les voitures n'en ont jamais besoin~; leurs
\omterm{def:g5:everyday-energy:energy}{énergies} meurent dans des freins ordinaires.
\end{solution}

\begin{solution}{pb:g9:kinetic-energy-safety:1}
\textbf{1.} $50 \div 3.6 = \qty{13.9}{m/s}$~;
$60 \div 3.6 = \qty{16.7}{m/s}$.
\textbf{2.} \qty{13.9}{m} et \qty{16.7}{m}.
\textbf{3.} \qty{13}{m}~; à soixante,
$13 \times (60/50)^2 = 13 \times 1.44 \approx \qty{18.7}{m}$.
\textbf{4.} Arrêt~: $26.9$ contre \qty{35.4}{m} --- le blanc de
l'affiche~: environ \emph{huit \omterm{def:g2:measuring-length:units}{mètres} et demi} de plus, deux longueurs
de voiture au-delà de l'endroit où le \omterm{def:g4:conductors-insulators:conductor}{conducteur} à cinquante s'est
arrêté.
\textbf{5.} À cinquante~: environ \qty{97}{kJ}. À quatre-vingt-dix
($v = \qty{25}{m/s}$)~:
$0.5 \times 1000 \times 625 = \qty{312}{kJ}$.
\textbf{6.} $97 \div 9.8 \approx \qty{10}{m}$~;
$312 \div 9.8 \approx \qty{32}{m}$.
\textbf{7.} Dix \omterm{def:g2:measuring-length:units}{mètres}~: une maison de trois étages~; trente-deux~: un
immeuble de dix étages. Proposition~: «~Un choc à 50, c'est une chute
du troisième étage. À 90, du dixième. Choisis ta fenêtre.~»
\textbf{8.} Les hauteurs mesurent honnêtement l'\omterm{def:g5:everyday-energy:energy}{énergie} qu'il faut
mettre à la retraite --- ce nombre-là, aucune ingénierie ne le change.
Les zones de déformation et les ceintures changent la \emph{manière}
de la retraite~: des \omterm{def:g2:measuring-length:units}{mètres} d'arrêt en douceur au lieu de \omterm{def:g2:measuring-length:units}{centimètres}
d'arrêt cruel --- et c'est pourquoi l'analogie de la chute exagère les
blessures dans une voiture moderne, et pourquoi l'\omterm{def:g5:everyday-energy:energy}{énergie} qu'elle
figure est réelle tout de même.
\textbf{9.} \qty{25}{m} à l'aveugle par seconde~; un coup d'œil de
deux secondes, \qty{50}{m} --- un demi-terrain avant même que ne
commence la seconde de réaction.
\textbf{10.} $65 + 25 = \qty{90}{m}$ --- le \omterm{def:g4:conductors-insulators:conductor}{conducteur} fatigué ajoute
une longueur d'autobus et demie au total du \omterm{def:g4:conductors-insulators:conductor}{conducteur} attentif.
\textbf{11.} Il a le plus tort aux \emph{basses} \omterm{def:g7:motion-average-speed:formula}{vitesses}~: là, la
part de freinage au carré est petite et la portion de réaction fait
l'essentiel de la distance d'arrêt --- en ville, la seconde
\emph{est} le danger. (À grande \omterm{def:g7:motion-average-speed:formula}{vitesse}, le freinage domine --- mais
\qty{25}{m} par seconde n'y sont pas non plus une broutille.)
\textbf{12.} Par exemple~: «~Dix de \omterm{def:g7:motion-average-speed:formula}{vitesse} en plus, dix \omterm{def:g2:measuring-length:units}{mètres}
d'arrêt en plus --- la \omterm{def:g7:motion-average-speed:formula}{vitesse} entre deux fois, dont une au carré. Un
choc est une chute~: l'\omterm{def:g5:everyday-energy:energy}{énergie} croît comme le carré de la \omterm{def:g7:motion-average-speed:formula}{vitesse}. Une
seconde de coup d'œil fait vingt-cinq \omterm{def:g2:measuring-length:units}{mètres} d'aveuglement ---
l'attention est la plus courte des \omterm{prop:g9:kinetic-energy-safety:stopping}{distances de freinage}.~»
\end{solution}
